\documentclass[aspectratio=169,10pt,fontset=fandol]{ctexbeamer}
\usepackage{amsmath,amssymb,mathtools,bm,mathrsfs}
\usepackage{graphicx,booktabs,array,tabularx,multirow}
\usepackage{tikz,pgfplots,xparse,xstring,mdframed}
\usetikzlibrary{arrows.meta,calc,angles,positioning,shapes.geometric,graphs}
\pgfplotsset{compat=1.18}

\setbeamerfont{normal text}{size=\normalsize}
\setbeamerfont{title}{size=\large}
\setbeamerfont{subtitle}{size=\normalsize}
\setbeamerfont{frametitle}{size=\large}
\setbeamerfont{framesubtitle}{size=\small}
\setbeamerfont{block title}{size=\normalsize}
\setbeamerfont{block body}{size=\normalsize}
\setbeamerfont{section in toc}{size=\normalsize}
\setbeamerfont{section in head/foot}{size=\footnotesize}
\setbeamerfont{author in head/foot}{size=\footnotesize}
\setbeamerfont{title in head/foot}{size=\footnotesize}
\setbeamerfont{date in head/foot}{size=\footnotesize}
\setbeamertemplate{navigation symbols}{}
\setbeamertemplate{itemize item}{\normalsize\textbullet}
\setbeamertemplate{itemize subitem}{\small\textbullet}
\setbeamercolor{structure}{fg=blue!65!black}
\setbeamercolor{block title}{fg=blue!65!black,bg=blue!8}
\setbeamercolor{block body}{fg=black,bg=blue!3}
\setbeamertemplate{footline}{\leavevmode\hbox to\paperwidth{\hfill{\footnotesize\insertframenumber}\hspace*{1em}}\vskip3pt}

% ctex fontset=fandol; system fonts are not redistributed.

\newenvironment{proofdetail}{\par\noindent\textbf{证明：}\quad}{\par}
\renewenvironment{solution}{\par\noindent\textbf{解：}\quad}{\par}
\newenvironment{custom}[1]{\par\textbf{#1}}{\par}
\title{数列}
\begin{document}
\section{图解}
% FRAME-ID F001
\typeout{MATH-AGENT-FRAME-BEGIN F001}
\begin{frame}[t,fragile]{目录}
\normalsize
{\color{blue}\normalsize 1. 图解\par}\vspace{0.22em}
\end{frame}
\typeout{MATH-AGENT-FRAME-END F001}
% FRAME-ID F002
\typeout{MATH-AGENT-FRAME-BEGIN F002}
\begin{frame}[t,fragile]{图解}
\normalsize
{\normalsize
% PRESENTATION-SCAFFOLD isotropic coordinates; node font size unchanged
\tikzset{every picture/.append style={scale=0.75}}
% SOURCE-BEGIN U0001 0 99
\begin{custom}{几何}\begin{center}\begin{tikzpicture}\draw (0,0)--(1,1);\end{tikzpicture}\end{center}
% SOURCE-END U0001
% PRESENTATION-SCAFFOLD close continued environment
\end{custom}
}
\end{frame}
\typeout{MATH-AGENT-FRAME-END F002}
% FRAME-ID F003
\typeout{MATH-AGENT-FRAME-BEGIN F003}
\begin{frame}[t,fragile]{图解（续1）}
\footnotesize
{\footnotesize
% REPEATED-CONTEXT environment continuation
\begin{custom}{几何}
% SOURCE-BEGIN U0001 99 646
\begin{itemize}\item 解释。解释。解释。解释。解释。解释。解释。解释。解释。解释。解释。解释。解释。解释。解释。解释。解释。解释。解释。解释。解释。解释。解释。解释。解释。解释。解释。解释。解释。解释。解释。解释。解释。解释。解释。解释。解释。解释。解释。解释。
\item 解释。解释。解释。解释。解释。解释。解释。解释。解释。解释。解释。解释。解释。解释。解释。解释。解释。解释。解释。解释。解释。解释。解释。解释。解释。解释。解释。解释。解释。解释。解释。解释。解释。解释。解释。解释。解释。解释。解释。解释。
\item 解释。解释。解释。解释。解释。解释。解释。解释。解释。解释。解释。解释。解释。解释。解释。解释。解释。解释。解释。解释。解释。解释。解释。解释。解释。解释。解释。解释。解释。解释。解释。解释。解释。解释。解释。解释。解释。解释。解释。解释。
\item 解释。解释。解释。解释。解释。解释。解释。解释。解释。解释。解释。解释。解释。解释。解释。解释。解释。解释。解释。解释。解释。解释。解释。解释。解释。解释。解释。解释。解释。解释。解释。解释。解释。解释。解释。解释。解释。解释。解释。解释。\end{itemize}\end{custom}
% SOURCE-END U0001
}
\end{frame}
\typeout{MATH-AGENT-FRAME-END F003}
\end{document}
